\chapter*{Lời cảm ơn}
\addcontentsline{toc}{chapter}{Lời cảm ơn}
Nhóm thực hiện xin trân trọng cảm ơn cô Nguyễn Thị Hồng Tú, giảng viên môn Nhận dạng mẫu, đã cung cấp nền tảng kiến thức về biểu diễn dữ liệu, phân tích thành phần chính, phân tích phân biệt Fisher và mô hình Markov ẩn. Các nội dung này là cơ sở để nhóm xây dựng phương án nhận diện chuỗi ký tự biển số và thiết kế phép đối sánh với máy vectơ hỗ trợ.

Nhóm xin cảm ơn Phân hiệu Trường Đại học Thủy Lợi và Bộ môn Công nghệ thông tin đã tạo điều kiện học tập, nghiên cứu và thực hành. Nhóm đồng thời ghi nhận đóng góp của các tác giả bộ dữ liệu VNLP, các công trình nghiên cứu về HOG, SVM, FLD và HMM, cùng cộng đồng phát triển OpenCV, NumPy, SciPy và scikit-learn.

Trong phạm vi đồ án, nhóm chú trọng tính nhất quán giữa mô hình đã huấn luyện, dữ liệu đánh giá và chương trình minh họa. Bộ dữ liệu ký tự nội bộ được sử dụng để bổ sung các lớp thiếu mẫu. Đặc điểm nhãn yếu và tính độc lập của tập kiểm thử được trình bày để xác định đúng phạm vi kết luận.

Nhóm mong nhận được ý kiến chuyên môn của giảng viên để tiếp tục hoàn thiện phương pháp nhận dạng và quy trình thực nghiệm.

\vspace{1cm}
\begin{flushright}
Thành phố Hồ Chí Minh, ngày 26 tháng 10 năm 2026\\[0.3cm]
Nhóm thực hiện
\end{flushright}

\chapter*{Nhận xét của giảng viên}
\addcontentsline{toc}{chapter}{Nhận xét của giảng viên}
\noindent Giảng viên: Nguyễn Thị Hồng Tú

\vspace{0.5cm}
\foreach \n in {1,...,14}{\noindent\rule{\textwidth}{0.2pt}\par\vspace{0.5cm}}
\begin{flushright}
Ngày \rule{0.8cm}{0.2pt} tháng \rule{0.8cm}{0.2pt} năm \rule{1cm}{0.2pt}\\[0.4cm]
Giảng viên nhận xét\\[1.3cm]
\end{flushright}
